\chapter{Passive Distributed Control}
\label{ch:passive-control}

\epigraph{%
  The most efficient movement is one where the body's own mechanics
  do most of the work, and the nervous system merely guides.
}{--- Adapted from Frans Bosch}

\section{The Frans Bosch Framework}

Frans Bosch's work on strength training and coordination \cite{Bosch2015}
provides a practitioner-oriented framework for understanding how elite athletes
organize movement. His key insights, translated into the language of this series:

\begin{insight}{Self-Organization in Complex Movement}{bosch-self-org}
  Bosch argues that complex movements are not centrally programmed but
  \textbf{self-organize} through the interaction of:
  \begin{enumerate}
    \item \textbf{Task constraints}: the goal shapes the solution space
    \item \textbf{Organismic constraints}: the body's anatomy and current state
    \item \textbf{Environmental constraints}: gravity, ground, implements
  \end{enumerate}
  The role of the nervous system is not to specify every muscle activation but to
  set the \emph{initial conditions} and \emph{constraints} that allow self-organization
  to produce the desired movement.

  In the language of Volume~I: the CNS designs the \emph{drift field} $f(\state)$
  of the system (through postural configuration and muscle tone) and lets the
  natural dynamics do the work.
\end{insight}

\section{Passive Dynamic Walking}

The most striking example of passive dynamics is passive dynamic walking
(McGeer, 1990): a simple mechanical biped with knees, feet, and hip (no motors)
walks down a gentle slope in stable limit cycles with zero active control. This
demonstrates that complex coordinated movement can emerge from mechanical design
and gravity alone.

\noindent
Mathematically, the walker is a system with zero control input: $\dot{\state} = f(\state)$
is purely drift dynamics with no control term $G(\state) \bm{u}$. The stability analysis
of periodic orbits in such systems is a central topic in Agrachev \& Sachkov \cite{AgrachevSachkov2004},
Chapter~2 (affine systems with zero control) and Chapter~4 (limit cycles and stability):
the existence of stable walking gaits relies on eigenvalue analysis of the Poincaré return map,
a classical dynamical systems tool that forms the foundation of their geometric approach.
The most striking example of passive dynamics is passive dynamic walking,
demonstrated by McGeer (1990) \cite{McGeer1990}: a simple mechanical biped
can walk down a gentle slope with \emph{no actuators and no control}:

\begin{lstlisting}[language=Python, caption={Passive dynamic walker simulation.}]
import numpy as np
from scipy.integrate import solve_ivp

def passive_walker_dynamics(
    t: float,
    state: np.ndarray,
    M: float = 10.0,
    m: float = 5.0,
    l: float = 1.0,
    g: float = 9.81,
    slope: float = 0.05
) -> np.ndarray:
    """Dynamics of a 2-DOF passive walker (compass gait).

    Parameters
    ----------
    t : float
        Time.
    state : np.ndarray
        [theta_stance, theta_swing, d_theta_stance, d_theta_swing]
    M : float
        Hip mass (kg).
    m : float
        Leg mass at foot (kg).
    g : float
        Gravitational acceleration (m/s^2).
    slope : float
        Ground slope (radians).

    Returns
    -------
    np.ndarray
        State derivatives.
    """
    th_st, th_sw, dth_st, dth_sw = state

    # Mass matrix
    M11 = (M + m) * l**2
    M12 = -m * l**2 * np.cos(th_st - th_sw)
    M21 = M12
    M22 = m * l**2

    # Gravity and Coriolis
    C1 = -m * l**2 * dth_sw**2 * np.sin(th_st - th_sw) \
         - (M + m) * g * l * np.sin(th_st - slope)
    C2 = m * l**2 * dth_st**2 * np.sin(th_st - th_sw) \
         + m * g * l * np.sin(th_sw - slope)

    # Solve M * ddq = C (NO control input!)
    det = M11 * M22 - M12 * M21
    ddth_st = (M22 * C1 - M12 * C2) / det
    ddth_sw = (M11 * C2 - M21 * C1) / det

    return np.array([dth_st, dth_sw, ddth_st, ddth_sw])


# Simulate one step of passive walking
state0 = np.array([0.2, -0.4, -0.5, 1.0])
sol = solve_ivp(passive_walker_dynamics, [0, 1.5], state0,
                max_step=0.001, events=None)
print(f"Passive walker: simulated {sol.t[-1]:.2f}s "
      f"with {len(sol.t)} timesteps")
print(f"Final stance angle: {sol.y[0, -1]:.3f} rad")
print("(No actuators, no controller --- just gravity and mechanics)")
\end{lstlisting}

\section{Passive Dynamics in the Golf Swing}

The golf swing exemplifies passive dynamics exploitation:
\begin{enumerate}
  \item \textbf{Backswing}: energy stored in elastic tissues (muscles, tendons, fascia)
  \item \textbf{Transition}: elastic recoil initiates the downswing passively
  \item \textbf{Wrist release}: centrifugal force passively uncocks the wrists
        (the double pendulum effect from Volume~II, Chapter~5)
  \item \textbf{Follow-through}: passive deceleration through tissue compliance
\end{enumerate}

In the ZTCF framework of Volume~I, Chapter~7: the drift contribution $f(\state)$
accounts for $\sim$60\% of the clubhead speed. The control contribution
$G(\state)\control$ provides the remaining $\sim$40\%.

Trajectory design exploiting drift is formalized in Agrachev \& Sachkov \cite{AgrachevSachkov2004},
Chapter~5 (optimal control with drift): when the drift term dominates the control term, optimal
trajectories naturally evolve along low-cost paths in the accessible set, allowing the system to
reach distant targets with minimal control effort. The golf swing is a biological instance of
this principle.

\section{Bernstein's Stage 3: Exploiting DOF}

The connection to Bernstein's learning framework (Chapter~1) is direct:
\begin{itemize}
  \item The expert athlete has reached Stage 3 --- exploiting degrees of freedom
  \item They have learned to configure the body so that passive dynamics
        \emph{automatically} produce the desired motion
  \item The control problem is simplified: instead of commanding every muscle,
        set up the initial conditions and let physics do the work
\end{itemize}

\section*{Exercises}

\begin{enumerate}
  \item \textbf{Passive Walking.} Using the provided code, find the set of
        initial conditions that produce a stable limit cycle (periodic gait).

  \item \textbf{Energy Analysis.} For the passive walker, plot kinetic and
        potential energy over one stride. Where does the energy come from?

  \item \textbf{Wrist Release.} Model the wrist release as a passive double
        pendulum. For what initial angular velocity of the arm does the club
        reach maximum speed at the bottom?

  \item \textbf{(Programming.)} Add a small actuator to the passive walker
        and design a minimal controller that stabilizes the gait on level ground.
\end{enumerate}
